% !TEX root = ../main.tex
%% Corrigés — Chapitre 14 — Représentations paramétriques et équations cartésiennes
%% Tous les résultats ont été vérifiés par programme (calcul exact).

\corrigesChapitre{14}{Représentations paramétriques}

\begin{deuxcolonnes}

\begin{corrige}{1}{Une droite du plan}
\begin{questions}
  \item $\res{d:\ x=3+t,\ y=-2+2t}$, $t\in\R$.
  \item $3+t=5$ donne $t=2$, et $-2+2\times2=2$ : $B\in d$.
  \item $t=-2$ : $\res{(1\pv-6)}$. Ordonnée $4$ : $-2+2t=4$, $t=3$, point $\res{(6\pv4)}$.
  \item Par exemple $\res{(2\pv4)}$, ou tout vecteur non nul colinéaire à $\vec u$.
\end{questions}
\end{corrige}

\begin{corrige}{2}{Une droite de l'espace}
\begin{questions}
  \item Une représentation paramétrique de $d$ est
    \[\res{\begin{cases}x=1+2t\\y=2-t\\z=-3+4t\end{cases}}\qquad t\in\R.\]
  \item Pour $B$ : $t=1$ convient pour les trois coordonnées, $\res{B\in d}$. Pour $C$ : la première équation donne $t=-\frac12$, et
    alors $y=\frac52\neq3$ : $\res{C\notin d}$.
  \item $t=2$ : $\res{(5\pv0\pv5)}$. Cote $5$ : $-3+4t=5$, $t=2$ : c'est le même point.
\end{questions}
\end{corrige}

\begin{corrige}{3}{Plusieurs représentations d'une même droite}
\begin{questions}
  \item $\res{d:\ x=2+t,\ y=-1+2t,\ z=3-t}$, $t\in\R$.
  \item $M$ : $t=2$, et $N$ : $t=1$. Les deux points sont sur $d$. Comme $\vect{AM}\,(2\pv4\pv-2)=2\vec u$, c'est un
    $\res{\text{vecteur directeur}}$ de $d$.
  \item $\res{P(0\pv-5\pv5)}$.
  \item $x=4+2k$, $y=3+4k$, $z=1-2k$, $k\in\R$. Le point $A$ correspond à $\res{k=-1}$.
\end{questions}
\end{corrige}

\begin{corrige}{4}{Du paramétrique au cartésien dans le plan}
\begin{questions}
  \item $t=\frac{x-1}{2}$ donne $y=3-2(x-1)$, soit $\res{2x+y-5=0}$ ; vecteur normal $\res{\vec n\,(2\pv1)}$.
  \item $\res{d':\ 2x+y=0}$.
  \item Les points $(2k\pv k)$ de la droite dirigée par $\vec n$ sont sur $d$ si $5k-5=0$ : $\res{H(2\pv1)}$. La distance vaut
    $OH=\res{\sqrt5}$.
\end{questions}
\end{corrige}

\begin{corrige}{5}{Équation cartésienne d'un plan}
\begin{questions}
  \item $3x-y+2z+d=0$ et $A$ donne $3-2-2+d=0$ : $\res{3x-y+2z+1=0}$. Pour $B$ : $6-5+2+1=4\neq0$ ; pour $O$ : $1\neq0$. Aucun
    des deux n'appartient à $\mathcal{P}$.
  \item $\vec n\cdot\vec{u_1}=3-3=0$ et $\vec n\cdot\vec{u_2}=-2+2=0$, avec $\vec{u_1}$, $\vec{u_2}$ non colinéaires :\\
    $x=1+s$, $y=2+3s+2t$, $z=-1+t$.
  \item $\res{3x-y+2z=0}$.
\end{questions}
\end{corrige}

\begin{corrige}{6}{D'une représentation à l'autre}
\begin{questions}
  \item $x=1+s$, $y=1+t$, $z=2s-t$. Avec $s=x-1$ et $t=y-1$ : $z=2(x-1)-(y-1)$, soit $\res{2x-y-z-1=0}$.
  \item $\vec n\cdot\vec u=2-2=0$ et $\vec n\cdot\vec v=-1+1=0$ : c'est cohérent, $\vec n$ est le vecteur normal lu sur l'équation.
  \item $6-2-3-1=0$ : $M\in\mathcal{P}$, avec $\res{s=2}$ et $\res{t=1}$.
  \item $\res{x=s,\ y=t,\ z=3-s-t}$, $(s\pv t)\in\R^2$.
\end{questions}
\end{corrige}

\begin{corrige}{7}{Quatre droites et leurs positions}
\begin{questions}
  \item $d_1$ : $(1+2t\pv-1+t\pv2-t)$. $d_2$, dirigée par $\frac13\vect{BC}\,(1\pv-1\pv1)$ : $(k\pv2-k\pv1+k)$. $d_4$ :
    $(-1+2\lambda\pv\lambda\pv3-\lambda)$. $d_3$ passe par $(3\pv0\pv1)$, dirigée par $(1\pv-2\pv1)$.
  \item Les deux premières équations donnent $t=\frac23$, $k=\frac73$, mais la troisième n'est pas vérifiée ; les vecteurs ne sont
    pas colinéaires : $d_1$ et $d_2$ sont $\res{\text{non coplanaires}}$.
  \item $t=1$ et $s=0$ conviennent : $\res{E(3\pv0\pv1)}$.
  \item $k=1$ et $\lambda=1$ conviennent : $\res{G(1\pv1\pv2)}$.
\end{questions}
\end{corrige}

\begin{corrige}{8}{Une droite et un plan}
\begin{questions}
  \item $\vec u\,(2\pv-1\pv1)$, $\vec n\,(1\pv2\pv-1)$ et $\vec u\cdot\vec n=2-2-1=\res{-1}$.
  \item $\vec u\cdot\vec n\neq0$ : $d$ n'est pas parallèle à $\mathcal{P}$. En reportant : $12-t=0$, d'où $t=12$ et
    $\res{I(25\pv-9\pv11)}$.
  \item $A+t\vec n$ dans l'équation : $12+6t=0$, $t=-2$ :\\ $\res{H(-1\pv-1\pv1)}$ et $AH=2\norme{\vec n}=\res{2\sqrt6}$.
\end{questions}
\end{corrige}

\begin{corrige}{9}{Deux droites sécantes}
\begin{questions}
  \item $\vec{u_1}\,(1\pv2\pv1)$ et $\vec{u_2}\,(1\pv1\pv1)$ ne sont pas colinéaires : $\res{\text{non parallèles}}$.
  \item $t=1+s$ et $2t=2+s$ donnent $s=0$, $t=1$, et la troisième équation est vérifiée : $\res{I(1\pv2\pv2)}$.
  \item $\cos\theta=\frac{4}{\sqrt6\sqrt3}=\frac{2\sqrt2}{3}\approx0{,}943$, soit $\theta\approx\res{19^\circ}$.
\end{questions}
\end{corrige}

\begin{corrige}{10}{Deux droites non coplanaires}
\begin{questions}
  \item $(2\pv-1\pv1)$ et $(1\pv2\pv3)$ ne sont pas colinéaires : $\res{\text{non parallèles}}$.
  \item $1+2t=2+s$ et $3-t=1+2s$ donnent $t=\frac45$, $s=\frac35$. Mais $-1+t=-\frac15$ et $3s=\frac95$ : la troisième équation est
    fausse.
  \item Ni parallèles ni sécantes : $\res{\text{non coplanaires}}$. Aucun plan ne les contient toutes les deux.
\end{questions}
\end{corrige}

\end{deuxcolonnes}

%% ---------------------------------------------------------------------
\corrigesSouples
\begin{corrige}{11}{Deux plans sécants}
\begin{questions}
  \item Les vecteurs normaux $\vec{n_1}\,(1\pv1\pv1)$ et $\vec{n_2}\,(2\pv-1\pv1)$ ne sont pas colinéaires : les plans sont
    $\res{\text{sécants}}$. Comme $\vec{n_1}\cdot\vec{n_2}=2\neq0$, ils ne sont pas perpendiculaires.
  \item Avec $y=t$ : $x+z=1-t$ et $2x+z=-3+t$. Par différence $x=-4+2t$, puis $z=5-3t$ :
    $\res{\Delta:\ (-4+2t\pv t\pv5-3t)}$, $t\in\R$.
  \item $-4+0+5-1=0$ et $-8-0+5+3=0$ : $M$ appartient aux deux plans (c'est le point de paramètre $0$). Le vecteur directeur
    $(2\pv1\pv-3)$ vérifie $2+1-3=0$ et $4-1-3=0$.
  \item $\cos\theta=\frac{2}{\sqrt3\sqrt6}=\frac{\sqrt2}{3}\approx0{,}471$, soit $\theta\approx\res{62^\circ}$.
\end{questions}
\end{corrige}

\begin{corrige}{12}{Synthèse : droites, plan et intersection}
\begin{questions}
  \item $\vect{AB}\,(2\pv1\pv-2)$ : $\res{(AB):\ (1+2t\pv t\pv2-2t)}$, $t\in\R$.
  \item $\vect{AC}\,(1\pv-1\pv1)$ : $\vec n\cdot\vect{AB}=2+4-6=0$ et $\vec n\cdot\vect{AC}=1-4+3=0$, avec $\vect{AB}$ et $\vect{AC}$ non
    colinéaires : $\vec n$ est normal à $(ABC)$. L'équation $x+4y+3z+d=0$ passe par $A$ si $d=-7$ : $\res{x+4y+3z-7=0}$.
  \item $0-4+12-7=1\neq0$ : $\res{D\notin(ABC)}$.
  \item Le vecteur directeur $(1\pv-1\pv2)$ de $d$ vérifie $\vec n\cdot(1\pv-1\pv2)=3\neq0$ : $d$ coupe le plan. En reportant :
    $(1+t)+4(2-t)+3(3+2t)-7=11+3t=0$, d'où $t=-\frac{11}3$ et $\res{I\left(-\tfrac83\pv\tfrac{17}{3}\pv-\tfrac{13}{3}\right)}$.
\end{questions}
\end{corrige}

\begin{corrige}{13}{Une pyramide à base rectangulaire}
\begin{questions}
  \item $\vect{AB}\,(3\pv0\pv0)=\vect{DC}$ : $ABCD$ est un parallélogramme ; $\vect{AB}\cdot\vect{AD}=0$ avec $\vect{AD}\,(0\pv2\pv0)$ : c'est un
    rectangle, d'aire $3\times2=\res{6}$.
  \item $\vect{SD}\,(-1\pv1\pv-3)$ : $\res{(SD):\ (1-t\pv1+t\pv3-3t)}$. $(SA)$ passe par $A$, dirigée par $\vect{AS}\,(1\pv1\pv3)$ :
    $\res{(SA):\ (u\pv u\pv3u)}$.
  \item $\vect{SP}=\frac14\vect{SD}$ donne $\res{P\left(\tfrac34\pv\tfrac54\pv\tfrac94\right)}$.
  \item $\vect{BP}\left(-\frac94\pv\frac54\pv\frac94\right)$ et $\vect{BC}\,(0\pv2\pv0)$ ne sont pas colinéaires. Les trois points vérifient
    $x+z-3=0$ : $\frac34+\frac94-3=0$, $3+0-3=0$ et $3+0-3=0$. C'est donc une équation du plan $(PBC)$.
  \item $u+3u-3=0$ donne $u=\frac34$ : $\res{Q\left(\tfrac34\pv\tfrac34\pv\tfrac94\right)}$. La section de la pyramide par $(PBC)$
    est le trapèze $BCPQ$, car $(PQ)\parallel(AD)\parallel(BC)$.
  \item $\Delta$ est dirigée par $\vec n\,(1\pv0\pv1)$ : $(1+t\pv1\pv3+t)$. Dans l'équation du plan : $1+t+3+t-3=0$, $t=-\frac12$ :
    $\res{H\left(\tfrac12\pv1\pv\tfrac52\right)}$, et $SH=\frac12\norme{\vec n}=\res{\tfrac{\sqrt2}{2}}$.
  \item $\vect{KP}\left(-\frac94\pv0\pv\frac94\right)$ : $PK=\frac94\sqrt2$. L'aire de $PBC$ vaut $\frac12\times BC\times PK=\res{\tfrac{9\sqrt2}{4}}$
    et le volume du tétraèdre $SPBC$ vaut $\frac13\times\frac{9\sqrt2}{4}\times\frac{\sqrt2}{2}=\res{\tfrac34}$.
\end{questions}
\end{corrige}

\begin{corrige}{14}{Un tétraèdre trirectangle}
\begin{questions}
  \item $6\times1-6=0$, $3\times2-6=0$ et $2\times3-6=0$ : les trois points vérifient l'équation, qui est donc celle du plan $(ABC)$
    (les points ne sont pas alignés).
  \item $\Delta$ est dirigée par le vecteur normal $\vec n\,(6\pv3\pv2)$ : $\res{\Delta:\ (6t\pv3t\pv2t)}$, $t\in\R$. En reportant :
    $36t+9t+4t-6=0$, soit $t=\frac{6}{49}$, d'où $\res{H\left(\tfrac{36}{49}\pv\tfrac{18}{49}\pv\tfrac{12}{49}\right)}$.
  \item $\norme{\vec n}=\sqrt{36+9+4}=7$, donc $OH=\frac{6}{49}\times7=\res{\tfrac67}$.
  \item Avec la base $OAB$, rectangle en $O$, et la hauteur $OC$ : $V=\frac13\times\frac{1\times2}{2}\times3=\res{1}$. Avec la base $ABC$
    et la hauteur $OH$ : $1=\frac13\times\mathcal{A}_{ABC}\times\frac67$, d'où $\mathcal{A}_{ABC}=\res{\tfrac72}$.
  \item $\vect{AH}\left(-\frac{13}{49}\pv\frac{18}{49}\pv\frac{12}{49}\right)$ et $\vect{BC}\,(0\pv-2\pv3)$ : $\frac{-36+36}{49}=0$.
    $\vect{BH}\left(\frac{36}{49}\pv-\frac{80}{49}\pv\frac{12}{49}\right)$ et $\vect{AC}\,(-1\pv0\pv3)$ : $\frac{-36+36}{49}=0$. Les droites
    $(AH)$ et $(BH)$ sont deux hauteurs du triangle $ABC$ : $H$ est son $\res{\text{orthocentre}}$.
\end{questions}
\end{corrige}
